% GENERATED FILE. Edit canonical lessons, then run npm run generate:books.
% canonical-source-hash: fnv1a64:76f2f72968465aff
% canonical-lessons: SA-C462-sathah, SA-C462-kutilah, SA-C462-apalapati, SA-C462-visvasaniyah, SA-C462-visvasaghatah

\chapter{Honest and False}
\label{ch:sa-honest-and-false}
\hlchaptermodality{462}

\begin{chapteropening}
\textbf{By the end of this chapter:} \emph{I can say whether someone can be trusted.}

Complete the last lesson of chapter 462: I can say whether someone can be trusted.
\end{chapteropening}

\section[\texorpdfstring{\sk{शठः} (śaṭhaḥ)}{शठः (śaṭhaḥ)}]{\sk{शठः} (śaṭhaḥ) — deceitful}

\label{lesson:SA-C462-sathah}



\begin{quote}
Before the new one: say the Sanskrit for heartless, then the Sanskrit for bears malice.
\end{quote}

\subsection*{You'll want to know: \sk{शठः}}
\textbf{\sk{शठः}} — \emph{śaṭhaḥ} — \textquotedblleft{}deceitful\textquotedblright{}.

\begin{quote}
\sk{शठः} \sk{वणिक्} \sk{मिथ्या} \sk{वदति।}
\end{quote}

\begin{quote}
\emph{śaṭhaḥ vaṇik mithyā vadati.}
\end{quote}

\begin{quote}
The deceitful merchant tells lies.
\end{quote}

\begin{cousinweb}
\textbf{\sk{शठः}} (\emph{śaṭhaḥ}) MW gives \textquotedblleft{}false, deceitful, fraudulent\textquotedblright{} and \textquotedblleft{}a cheat, a rogue\textquotedblright{}; MW gives no derivation, so its origin is uncertain (MW).
\end{cousinweb}

\begin{grammarlens}[title={What you've built}]
One more word for describing what you lived and felt.
\end{grammarlens}

\subsection*{Your turn}
\begin{itemize}
\raggedright
  \item \emph{Say it:} \emph{śaṭhaḥ}
  \item \emph{Say it:} \emph{śaṭhaḥ}, once more
  \item \emph{Say it:} \emph{śaṭhaḥ vaṇik mithyā vadati.} — The deceitful merchant tells lies.
  \item \emph{Recall:} say the Sanskrit for heartless, then the Sanskrit for bears malice, then say \emph{śaṭhaḥ} again
\end{itemize}

\begin{tcolorbox}[breakable,colback=teal!4,colframe=teal!35!black,title={Before you move on}]
What does \sk{शठः} mean? (\textquotedblleft{}Deceitful\textquotedblright{}.) Say it once more.
\end{tcolorbox}
\section[\texorpdfstring{\sk{कुटिलः} (kuṭilaḥ)}{कुटिलः (kuṭilaḥ)}]{\sk{कुटिलः} (kuṭilaḥ) — crooked; devious}

\label{lesson:SA-C462-kutilah}



\begin{quote}
Before the new one: say the Sanskrit for bears malice, then the Sanskrit for deceitful.
\end{quote}

\subsection*{You'll want to know: \sk{कुटिलः}}
\textbf{\sk{कुटिलः}} — \emph{kuṭilaḥ} — \textquotedblleft{}crooked; devious\textquotedblright{}.

\begin{quote}
\sk{कुटिलः} \sk{जनः} \sk{सरलं} \sk{मित्रं} \sk{वञ्चयति।}
\end{quote}

\begin{quote}
\emph{kuṭilaḥ janaḥ saralaṃ mitraṃ vañcayati.}
\end{quote}

\begin{quote}
A devious man deceives his straightforward friend.
\end{quote}

\begin{cousinweb}
\textbf{\sk{कुटिलः}} (\emph{kuṭilaḥ}) MW's first sense is \textquotedblleft{}bent, crooked, curved\textquotedblright{}, and then \textquotedblleft{}dishonest, fraudulent\textquotedblright{}; MW gives no derivation, so its origin is uncertain (MW).
\end{cousinweb}

\begin{grammarlens}[title={What you've built}]
One more word for describing what you lived and felt.
\end{grammarlens}

\subsection*{Your turn}
\begin{itemize}
\raggedright
  \item \emph{Say it:} \emph{kuṭilaḥ}
  \item \emph{Say it:} \emph{kuṭilaḥ}, once more
  \item \emph{Say it:} \emph{kuṭilaḥ janaḥ saralaṃ mitraṃ vañcayati.} — A devious man deceives his straightforward friend.
  \item \emph{Recall:} say the Sanskrit for bears malice, then the Sanskrit for deceitful, then say \emph{kuṭilaḥ} again
\end{itemize}

\begin{tcolorbox}[breakable,colback=teal!4,colframe=teal!35!black,title={Before you move on}]
What does \sk{कुटिलः} mean? (\textquotedblleft{}Crooked; devious\textquotedblright{}.) Say it once more.
\end{tcolorbox}
\section[\texorpdfstring{\sk{अपलपति} (apalapati)}{अपलपति (apalapati)}]{\sk{अपलपति} (apalapati) — explains away; conceals}

\label{lesson:SA-C462-apalapati}



\begin{quote}
Before the new one: say the Sanskrit for deceitful, then the Sanskrit for devious.
\end{quote}

\subsection*{You'll want to know: \sk{अपलपति}}
\textbf{\sk{अपलपति}} — \emph{apalapati} — \textquotedblleft{}explains away; conceals\textquotedblright{}.

\begin{quote}
\sk{चौरः} \sk{सत्यम्} \sk{अपलपति।}
\end{quote}

\begin{quote}
\emph{cauraḥ satyam apalapati.}
\end{quote}

\begin{quote}
The thief explains away the truth.
\end{quote}

Past (he …): \sk{सः} \sk{सत्यम्} \sk{अपलपितवान्}, \emph{saḥ satyam apalapitavān}.

\begin{cousinweb}
\textbf{\sk{अपलपति}} (\emph{apalapati}) is built from \emph{apa}, \textquotedblleft{}away\textquotedblright{}, and the root √lap, \textquotedblleft{}to talk\textquotedblright{}: to talk a thing away; MW gives \textquotedblleft{}to explain away, deny, conceal\textquotedblright{} (MW).
\end{cousinweb}

\begin{grammarlens}[title={What you've built}]
One more word for describing what you lived and felt.
\end{grammarlens}

\subsection*{Your turn}
\begin{itemize}
\raggedright
  \item \emph{Say it:} \emph{apalapati}
  \item \emph{Say it:} \emph{apalapati}, once more
  \item \emph{Say it:} \emph{cauraḥ satyam apalapati.} — The thief explains away the truth.
  \item \emph{Recall:} say the Sanskrit for deceitful, then the Sanskrit for devious, then say \emph{apalapati} again
\end{itemize}

\begin{tcolorbox}[breakable,colback=teal!4,colframe=teal!35!black,title={Before you move on}]
What does \sk{अपलपति} mean? (\textquotedblleft{}Explains away; conceals\textquotedblright{}.) Say it once more.
\end{tcolorbox}
\section[\texorpdfstring{\sk{विश्वसनीयः} (viśvasanīyaḥ)}{विश्वसनीयः (viśvasanīyaḥ)}]{\sk{विश्वसनीयः} (viśvasanīyaḥ) — trustworthy; reliable}

\label{lesson:SA-C462-visvasaniyah}



\begin{quote}
Before the new one: say the Sanskrit for devious, then the Sanskrit for explains away.
\end{quote}

\subsection*{You'll want to know: \sk{विश्वसनीयः}}
\textbf{\sk{विश्वसनीयः}} — \emph{viśvasanīyaḥ} — \textquotedblleft{}trustworthy; reliable\textquotedblright{}.

\begin{quote}
\sk{विश्वसनीयः} \sk{वणिक्} \sk{उचितं} \sk{मूल्यं} \sk{वदति।}
\end{quote}

\begin{quote}
\emph{viśvasanīyaḥ vaṇik ucitaṃ mūlyaṃ vadati.}
\end{quote}

\begin{quote}
A trustworthy merchant names a fair price.
\end{quote}

\begin{cousinweb}
\textbf{\sk{विश्वसनीयः}} (\emph{viśvasanīyaḥ}) is built on \emph{vi-√śvas}, \textquotedblleft{}to trust\textquotedblright{}, the root of \textbf{\sk{विश्वसिति}}, with the ending \emph{-anīya}, \textquotedblleft{}to be ...-ed\textquotedblright{}; MW gives \textquotedblleft{}to be trusted or relied on, reliable, trustworthy\textquotedblright{} (MW).
\end{cousinweb}

\begin{grammarlens}[title={What you've built}]
One more word for describing what you lived and felt.
\end{grammarlens}

\subsection*{Your turn}
\begin{itemize}
\raggedright
  \item \emph{Say it:} \emph{viśvasanīyaḥ}
  \item \emph{Say it:} \emph{viśvasanīyaḥ}, once more
  \item \emph{Say it:} \emph{viśvasanīyaḥ vaṇik ucitaṃ mūlyaṃ vadati.} — A trustworthy merchant names a fair price.
  \item \emph{Recall:} say the Sanskrit for devious, then the Sanskrit for explains away, then say \emph{viśvasanīyaḥ} again
\end{itemize}

\begin{tcolorbox}[breakable,colback=teal!4,colframe=teal!35!black,title={Before you move on}]
What does \sk{विश्वसनीयः} mean? (\textquotedblleft{}Trustworthy; reliable\textquotedblright{}.) Say it once more.
\end{tcolorbox}
\section[\texorpdfstring{\sk{विश्वासघातः} (viśvāsaghātaḥ)}{विश्वासघातः (viśvāsaghātaḥ)}]{\sk{विश्वासघातः} (viśvāsaghātaḥ) — a betrayal of trust}

\label{lesson:SA-C462-visvasaghatah}



\begin{quote}
Before the new one: say the Sanskrit for explains away, then the Sanskrit for trustworthy.
\end{quote}

\subsection*{You'll want to know: \sk{विश्वासघातः}}
\textbf{\sk{विश्वासघातः}} — \emph{viśvāsaghātaḥ} — \textquotedblleft{}a betrayal of trust\textquotedblright{}.

\begin{quote}
\sk{मित्रस्य} \sk{विश्वासघातः} \sk{महत्} \sk{दुःखं} \sk{ददाति।}
\end{quote}

\begin{quote}
\emph{mitrasya viśvāsaghātaḥ mahat duḥkhaṃ dadāti.}
\end{quote}

\begin{quote}
A friend's betrayal brings great sorrow.
\end{quote}

A masculine noun.

\begin{cousinweb}
\textbf{\sk{विश्वासघातः}} (\emph{viśvāsaghātaḥ}) is built from \textbf{\sk{विश्वासः}}, \textquotedblleft{}trust\textquotedblright{}, and \emph{ghāta}, \textquotedblleft{}a blow, a killing\textquotedblright{}, from the root √han; MW gives \textquotedblleft{}violation of trust, treachery\textquotedblright{} (MW).
\end{cousinweb}

\begin{grammarlens}[title={What you've built}]
One more word for describing what you lived and felt.
\end{grammarlens}

\subsection*{Your turn}
\begin{itemize}
\raggedright
  \item \emph{Say it:} \emph{viśvāsaghātaḥ}
  \item \emph{Say it:} \emph{viśvāsaghātaḥ}, once more
  \item \emph{Say it:} \emph{mitrasya viśvāsaghātaḥ mahat duḥkhaṃ dadāti.} — A friend's betrayal brings great sorrow.
  \item \emph{Recall:} say the Sanskrit for explains away, then the Sanskrit for trustworthy, then say \emph{viśvāsaghātaḥ} again
\end{itemize}

\begin{tcolorbox}[breakable,colback=teal!4,colframe=teal!35!black,title={Before you move on}]
What does \sk{विश्वासघातः} mean? (\textquotedblleft{}A betrayal of trust\textquotedblright{}.) Say it once more.
\end{tcolorbox}
